% GENERATED FILE. Edit canonical lessons, then run npm run generate:books.
% canonical-source-hash: fnv1a64:c27eb1106650ffa4
% canonical-lessons: SA-C216-kusalah, SA-C216-nipunah, SA-C216-ajnah, SA-C216-vijnah, SA-C216-buddhiman

\chapter{Able, Expert, Ignorant, Wise, Intelligent}
\label{ch:sa-able-expert-ignorant-wise-intelligent}
\hlchaptermodality{216}

\begin{chapteropening}
\textbf{By the end of this chapter:} \emph{I can say able, expert, ignorant, wise and intelligent.}

Complete the last lesson of chapter 216: I can say able, expert, ignorant, wise and intelligent.
\end{chapteropening}

\section[\texorpdfstring{\sk{कुशलः} (kuśalaḥ)}{कुशलः (kuśalaḥ)}]{\sk{कुशलः} (kuśalaḥ) — able, well}

\label{lesson:SA-C216-kusalah}



\begin{quote}
Before the new one: say the Sanskrit for blunt, then the Sanskrit for valuable.
\end{quote}

\subsection*{You'll want to know: \sk{कुशलः}}
\textbf{\sk{कुशलः}} — \emph{kuśalaḥ} — \textquotedblleft{}able, well\textquotedblright{}.

\begin{grammarlens}[title={What you've built}]
One more word for this chapter.
\end{grammarlens}

\subsection*{Your turn}
\begin{itemize}
\raggedright
  \item \emph{Say it:} \emph{kuśalaḥ}
  \item \emph{Say it:} \emph{kuśalaḥ}, once more
  \item \emph{Say it:} say \emph{kuśalaḥ}
  \item \emph{Recall:} say the Sanskrit for blunt, then the Sanskrit for valuable, then say \emph{kuśalaḥ} again
\end{itemize}

\begin{tcolorbox}[breakable,colback=teal!4,colframe=teal!35!black,title={Before you move on}]
What does \sk{कुशलः} mean? (\textquotedblleft{}Able, well\textquotedblright{}.) Say it once more.
\end{tcolorbox}
\section[\texorpdfstring{\sk{निपुणः} (nipuṇaḥ)}{निपुणः (nipuṇaḥ)}]{\sk{निपुणः} (nipuṇaḥ) — expert}

\label{lesson:SA-C216-nipunah}



\begin{quote}
Before the new one: say the Sanskrit for valuable, then the Sanskrit for able.
\end{quote}

\subsection*{You'll want to know: \sk{निपुणः}}
\textbf{\sk{निपुणः}} — \emph{nipuṇaḥ} — \textquotedblleft{}expert\textquotedblright{}.

\begin{grammarlens}[title={What you've built}]
One more word for this chapter.
\end{grammarlens}

\subsection*{Your turn}
\begin{itemize}
\raggedright
  \item \emph{Say it:} \emph{nipuṇaḥ}
  \item \emph{Say it:} \emph{nipuṇaḥ}, once more
  \item \emph{Say it:} say \emph{nipuṇaḥ}
  \item \emph{Recall:} say the Sanskrit for valuable, then the Sanskrit for able, then say \emph{nipuṇaḥ} again
\end{itemize}

\begin{tcolorbox}[breakable,colback=teal!4,colframe=teal!35!black,title={Before you move on}]
What does \sk{निपुणः} mean? (\textquotedblleft{}Expert\textquotedblright{}.) Say it once more.
\end{tcolorbox}
\section[\texorpdfstring{\sk{अज्ञः} (ajñaḥ)}{अज्ञः (ajñaḥ)}]{\sk{अज्ञः} (ajñaḥ) — ignorant}

\label{lesson:SA-C216-ajnah}



\begin{quote}
Before the new one: say the Sanskrit for able, then the Sanskrit for expert.
\end{quote}

\subsection*{You'll want to know: \sk{अज्ञः}}
\textbf{\sk{अज्ञः}} — \emph{ajñaḥ} — \textquotedblleft{}ignorant\textquotedblright{}.

\begin{grammarlens}[title={What you've built}]
One more word for this chapter.
\end{grammarlens}

\subsection*{Your turn}
\begin{itemize}
\raggedright
  \item \emph{Say it:} \emph{ajñaḥ}
  \item \emph{Say it:} \emph{ajñaḥ}, once more
  \item \emph{Say it:} say \emph{ajñaḥ}
  \item \emph{Recall:} say the Sanskrit for able, then the Sanskrit for expert, then say \emph{ajñaḥ} again
\end{itemize}

\begin{tcolorbox}[breakable,colback=teal!4,colframe=teal!35!black,title={Before you move on}]
What does \sk{अज्ञः} mean? (\textquotedblleft{}Ignorant\textquotedblright{}.) Say it once more.
\end{tcolorbox}
\section[\texorpdfstring{\sk{विज्ञः} (vijñaḥ)}{विज्ञः (vijñaḥ)}]{\sk{विज्ञः} (vijñaḥ) — wise, knowing}

\label{lesson:SA-C216-vijnah}



\begin{quote}
Before the new one: say the Sanskrit for expert, then the Sanskrit for ignorant.
\end{quote}

\subsection*{You'll want to know: \sk{विज्ञः}}
\textbf{\sk{विज्ञः}} — \emph{vijñaḥ} — \textquotedblleft{}wise, knowing\textquotedblright{}.

\begin{grammarlens}[title={What you've built}]
One more word for this chapter.
\end{grammarlens}

\subsection*{Your turn}
\begin{itemize}
\raggedright
  \item \emph{Say it:} \emph{vijñaḥ}
  \item \emph{Say it:} \emph{vijñaḥ}, once more
  \item \emph{Say it:} say \emph{vijñaḥ}
  \item \emph{Recall:} say the Sanskrit for expert, then the Sanskrit for ignorant, then say \emph{vijñaḥ} again
\end{itemize}

\begin{tcolorbox}[breakable,colback=teal!4,colframe=teal!35!black,title={Before you move on}]
What does \sk{विज्ञः} mean? (\textquotedblleft{}Wise, knowing\textquotedblright{}.) Say it once more.
\end{tcolorbox}
\section[\texorpdfstring{\sk{बुद्धिमान्} (buddhimān)}{बुद्धिमान् (buddhimān)}]{\sk{बुद्धिमान्} (buddhimān) — intelligent}

\label{lesson:SA-C216-buddhiman}



\begin{quote}
Before the new one: say the Sanskrit for ignorant, then the Sanskrit for wise.
\end{quote}

\subsection*{You'll want to know: \sk{बुद्धिमान्}}
\textbf{\sk{बुद्धिमान्}} — \emph{buddhimān} — \textquotedblleft{}intelligent\textquotedblright{}.

\begin{grammarlens}[title={What you've built}]
One more word for this chapter.
\end{grammarlens}

\subsection*{Your turn}
\begin{itemize}
\raggedright
  \item \emph{Say it:} \emph{buddhimān}
  \item \emph{Say it:} \emph{buddhimān}, once more
  \item \emph{Say it:} say \emph{buddhimān}
  \item \emph{Recall:} say the Sanskrit for ignorant, then the Sanskrit for wise, then say \emph{buddhimān} again
\end{itemize}

\begin{tcolorbox}[breakable,colback=teal!4,colframe=teal!35!black,title={Before you move on}]
What does \sk{बुद्धिमान्} mean? (\textquotedblleft{}Intelligent\textquotedblright{}.) Say it once more.
\end{tcolorbox}
